% =====================================================================
%  PROTEUS — REVISION DROP-IN BLOCKS  (top-two fixes; PLOS Comput Biol)
%  The manuscript source (proteus.tex) is not in the repo — only the
%  compiled PDF. These are surgical, section-keyed drop-in blocks for the
%  author to paste in. Every number is generated by a committed script
%  (see REVISION_NOTES.md) and filled from the analysis JSON by
%  revision/fill_tokens.py — nothing here is hand-typed.
%
%  Double-brace tokens are substituted from data/processed/*.json.
%  After substitution this file is written to manuscript_revisions.filled.tex.
% =====================================================================


% ---------------------------------------------------------------------
% [BLOCK 0]  BYLINE / AFFILIATION  (replaces "Avery Karlin / AveryCorp
%            Institute, New York"). Matches the projGauge paper byline.
% ---------------------------------------------------------------------
\author{Avery Karlin\\[2pt]
\small The Annealing Signet Institute\\[3pt]
\small \textit{Avery Karlin is a graduate student at Columbia University,}\\[1pt]
\small \textit{Department of Applied Physics and Mathematics, and the}\\[1pt]
\small \textit{University of Colorado Boulder, Department of Computer Science.}\\[3pt]
\small ORCID: \href{https://orcid.org/0000-0003-3848-6782}{0000-0003-3848-6782}}


% ---------------------------------------------------------------------
% [BLOCK 1]  ABSTRACT  (full replacement). Reframes as a decision-relevant
%            cautionary benchmark; replaces "statistically indistinguishable
%            (p=0.30)" with the powered non-superiority result; folds in the
%            pLDDT-confound result; softens the strongest generalization.
% ---------------------------------------------------------------------
Protein structure prediction and structural search are increasingly proposed as a
route to discover divergent enzymes hidden in the metagenomic ``dark proteome'' ---
sequences beyond the reach of homology-based search --- with poly(ethylene
terephthalate) (PET) hydrolases a flagship target. Whether structure-first search
genuinely extends \emph{functional} discovery past the boundary of sequence homology
has rarely been tested directly. We built a calibrated structure-first pipeline ---
fold-class structural search, a catalytic Ser--His--Asp geometry gate, and a
size-invariant active-site-exposure score --- validated to cleanly separate known
PET hydrolases from non-PET serine hydrolases on crystal structures, and applied it
to the ESM Metagenomic Atlas ($\sim$37M predicted structures) with a properly
sampled random baseline and a per-anchor decomposition. Fold-class search enriches
serine hydrolases $\sim${{triad_ratio}}-fold over random, but the active-site-exposure
score is not selective: it admits $\sim${{base_rate_pct}}\% of random triad-bearing
hydrolases. Against an \emph{enlarged} random baseline ({{base_n}} triad-bearers),
the PET-hydrolase branch is \emph{not} enriched above this floor: a two one-sided
test at a pre-registered $\pm${{delta_pp}}-point margin places the largest plausible
enrichment (one-sided 95\% bound) at {{upper_bound_pp}} points
({{tost_descriptor}}). A maximum-sensitivity re-search retrieves more than nine
million divergent ($<$20\% identity) structural neighbors, yet these are \emph{not}
lower in predicted confidence than the baseline (median pLDDT {{div_med_pct}} vs.\
{{base_med_pct}}; the depletion is therefore not an artifact of model quality), and
they are \emph{depleted} in the exposed-cleft phenotype; molecular docking does not
discriminate PET hydrolases from negatives either. The only above-baseline signal is
confined to near-homolog sequence identities, where homology search already reaches.
We conclude that, \emph{for this operationalization of PET-hydrolyzing function},
structure-first triage is an excellent serine-hydrolase \emph{finder} but not a
PET-function \emph{discriminator}, and does not extend PET-hydrolase discovery beyond
sequence homology's reach; structure and sequence fail together in the divergent
tail. We release the full pipeline, the enlarged baseline, and intermediate data as a
reproducible, decision-relevant benchmark for claims of structural ``discovery''
beyond homology.


% ---------------------------------------------------------------------
% [BLOCK 2]  RESULTS Sec 2.3 — the central sentence(s). Replace the
%            "indistinguishable from the 42.9% random baseline (Fisher
%            p=0.30, rate ratio 0.76)" clause with the powered result.
% ---------------------------------------------------------------------
% REPLACE the sentence ending "...clears the exposure line at 32.65% (96/294
% triad-bearers; CI 27.6--38.2%) --- statistically indistinguishable from the 42.9%
% random baseline (Fisher p = 0.30, rate ratio 0.76)." WITH:

Screening the top 300 hits per anchor, the PET-hydrolase branch clears the exposure
line at 32.65\% (96/294 triad-bearers; Wilson 95\% CI 27.6--38.2\%). To test the
load-bearing claim --- that this rate is \emph{not enriched above} the random
serine-hydrolase floor --- we enlarged the random triad-bearing baseline from the
original $n=28$ to {{base_n}} triad-bearers ({{pooled_screened}} random Atlas
proteins drawn under an explicit new seed and piped through the unchanged pipeline at
the same pinned line; \nameref{sec:methods}). On this enlarged baseline the floor
rate is {{base_rate_pct}}\% ({{base_k}}/{{base_n}}; Wilson 95\% CI
{{base_wilson_lo_pct}}--{{base_wilson_hi_pct}}\%). The PET-branch rate is
{{diff_pp}} points relative to it (90\% CI {{ci90_lo_pp}}, {{ci90_hi_pp}}); a two
one-sided test at a pre-registered $\pm${{delta_pp}}-point equivalence margin
rejects enrichment above the floor (one-sided $p={{p_upper}}$), with the largest
plausible enrichment (one-sided 95\% upper bound) at {{upper_bound_pp}} points.
{{tost_sentence}} (For reference, the small original baseline gave 42.9\% [12/28] and
a non-significant Fisher $p=0.30$; that ``indistinguishable'' verdict reflected the
$n=28$ baseline's low power rather than evidence of equality.) The
acetylcholinesterase branch is lower still (8.75\%, 26/297).


% ---------------------------------------------------------------------
% [BLOCK 3]  RESULTS — new short paragraph for Sec 2.5 (the pLDDT confound),
%            inserted after the "depleted, not absent" sentence; cites the
%            new supplementary figure.
% ---------------------------------------------------------------------
Could the divergent tail's exposed-cleft depletion be an ESMFold prediction-error
artifact rather than biology? We tested this directly (the comparison the analysis
previously only flagged). Across the three bands --- the sub-25\% divergent screened
tail ($n={{div_n}}$), the enlarged random-baseline triad-bearers ($n={{base_n}}$),
and the near-homolog above-line hits ($n={{nearhom_n}}$) --- predicted confidence is
not lower in the divergent tail. Its global pLDDT (median {{div_med_pct}}) is
comparable to, and if anything slightly higher than, the random baseline (median
{{base_med_pct}}; Mann--Whitney $p={{plddt_mwu_p}}$, Cliff's $\delta={{plddt_cliffs}}$,
i.e.\ the divergent band is \emph{not} the lower-confidence group), with the
near-homologs highest (median {{nearhom_med_pct}}); the most divergent $<$20\%-identity
subset is likewise not depleted in pLDDT (median {{div_sub20_med_pct}}). A
pLDDT-matched re-screen --- subsampling the divergent triad-bearers to the baseline's
confidence distribution --- leaves the exposed-cleft rate essentially unchanged
({{matched_raw_pct}}\% raw vs.\ {{matched_pct}}\% matched), far below the
$\sim${{base_rate_pct}}\% floor. Prediction quality therefore does not explain the
depletion: it is a property of the structures the model is confident about
(\nameref{figS_plddt}).


% ---------------------------------------------------------------------
% [BLOCK 4]  DISCUSSION — significance reframing (replace/extend the
%            "These results are complementary..." paragraph) + soften.
% ---------------------------------------------------------------------
These results are best read not as a methodological caveat but as a decision-relevant,
reproducible benchmark for a question the field is actively betting on. The actionable
boundary is sharp: structure-first triage is an excellent serine-hydrolase
\emph{finder} --- the fold/triad stage enriches hydrolases $\sim${{triad_ratio}}-fold
over random --- but, with these descriptors, it is not a PET-function
\emph{discriminator}, and the sequence-identity reach at which any PET-specific signal
exists is quantified rather than assumed. Reports of metagenome-mined PET hydrolases
are complementary to, not in conflict with, this picture: such efforts typically
validate candidates that are themselves near-homologs of known enzymes, or that pass
explicit homology-seeded sequence filters; our analysis quantifies the boundary
within which those pipelines operate. The practical implication is cautionary: with
these descriptors, in-silico structure-first search should not be expected, on its
own, to deliver \emph{novel} PET hydrolases from the divergent dark proteome, and
claims of structural ``discovery'' beyond homology warrant explicit demonstration
against a random baseline and a sequence-identity reach analysis of the kind reported
here.


% ---------------------------------------------------------------------
% [BLOCK 5]  DISCUSSION — new prior-art paragraph (insert after BLOCK 4).
%            VERIFIED citations only (Crossref-checked); see References block.
% ---------------------------------------------------------------------
Our finding converges with an independent result from large-scale phylogenomics:
\citet{mutti2025} report that newly developed structure-based methods do \emph{not}
outperform standard sequence-based methods for tree inference, and trace the failure
to the same root cause we observe --- poorly predicted, low-confidence structures in
the regime where structure would otherwise add value. Proteus differs in target
(function \emph{discrimination} of a specific chemistry, not homology detection or
phylogeny; PET hydrolysis as the flagship test) and in design (an explicit random
baseline plus an identity-reach decomposition), but the conclusion rhymes: where
sequence signal is exhausted, predicted structure does not independently recover what
sequence misses. This is consistent rather than in tension with structure-first
mining \emph{successes}. Foldseek-based metagenomic enzyme mining recovers hundreds of
candidates --- for example, \citet{jaito2026} retrieve 711 putative novel lipases ---
but do so by following structural search with rigorous sequence-similarity filtering,
and machine-learning-guided PET-hydrolase discovery \citep{nortonbaker2025} likewise
surveys natural sequence diversity (here yielding dozens of validated PET hydrolases).
Those candidates are homology-reachable or sequence-filtered by construction, exactly
the regime our near-homolog gradient occupies; they corroborate, rather than
contradict, a homology-bounded reach for the structural signal itself.


% ---------------------------------------------------------------------
% [BLOCK 6]  LIMITATIONS — replace the pLDDT hand-wave. The original read:
%   "...prediction error at low identity cannot be excluded entirely and merits
%    an explicit pLDDT comparison between the divergent band and the baseline."
% ---------------------------------------------------------------------
% REPLACE that clause WITH:

The Atlas structures are ESMFold predictions; we mitigated model-quality confounds by
restricting to the high-confidence clustered set (pLDDT $>$ 0.7) and by coverage-gating
sequence-identity estimates, and we tested the residual prediction-error confound
directly: the divergent band is not lower in pLDDT than the random baseline (median
{{div_med_pct}} vs.\ {{base_med_pct}}; Cliff's $\delta={{plddt_cliffs}}$) and a
pLDDT-matched re-screen leaves the exposed-cleft depletion intact, so low-identity
model noise does not account for the result.


% ---------------------------------------------------------------------
% [BLOCK 7]  METHODS — two new sub-paragraphs (append under Materials and
%            methods, after "Baseline sampling").
% ---------------------------------------------------------------------
\textbf{Enlarged baseline and equivalence test.} To power the non-enrichment claim,
the random triad-bearing baseline was enlarged from $n=28$ to {{base_n}} triad-bearers
by drawing {{new_screened}} additional uniform-random Atlas proteins (seeded
byte-offset sampling over the HQ-clust30 representative list under a new explicit seed,
{{enlarge_seed}}; block $k$ drawn under seed {{enlarge_seed}}$+k$ for reproducibility),
de-duplicated against the original draw and piped through the unchanged S4$\to$S5
pipeline at the pinned line ($-1.1587$); pooled with the original 1,500-draw floor this
gives {{pooled_screened}} screened proteins. The non-enrichment claim was tested on the
difference $d$ between the PET-branch rate (96/294) and the baseline rate by two
one-sided tests (TOST) at a pre-registered equivalence margin $\delta=
{{delta_pp}}$~percentage points (a single configuration constant): equivalence is
declared only if the 90\% confidence interval of $d$ lies within $\pm\delta$;
non-superiority (the load-bearing direction) if the one-sided test of $H_0\!: d \ge
\delta$ is rejected. We report both one-sided $p$-values, the 90\% CI of $d$, and the
one-sided 95\% upper bound on $d$ (the largest plausible enrichment), alongside the
legacy two-proportion and Fisher tests.

\textbf{pLDDT confound test.} Global per-model pLDDT (mean CA B-factor) and an
active-site-local pLDDT (mean over the catalytic-triad residues and their immediate
neighbours) were computed from the cached ESMFold V0 models for the divergent
($<$25\% identity) screened band, the enlarged random-baseline triad-bearers, and the
near-homolog above-line hits. Distributions were compared with Mann--Whitney $U$ and
Kolmogorov--Smirnov tests and Cliff's $\delta$ effect size. Because the divergent band
was not lower in pLDDT, a pLDDT-matched re-screen subsampled the divergent
triad-bearers to the baseline's global-pLDDT histogram (seed 1729) and recomputed the
exposed-cleft rate.


% ---------------------------------------------------------------------
% [BLOCK 8]  DECLARATIONS — add Competing Interests + Funding (PLOS requires
%            both). Keep the existing Gen-AI disclosure and Data availability.
% ---------------------------------------------------------------------
\section*{Competing interests}
The author has declared that no competing interests exist.

\section*{Funding}
The author received no specific funding for this work. The Annealing Signet Institute
provided no targeted financial support; compute was self-funded. The funders had no
role in study design, data collection and analysis, decision to publish, or
preparation of the manuscript.


% ---------------------------------------------------------------------
% [BLOCK 9]  REFERENCES — three NEW entries (Crossref-verified). Add to the
%            bibliography; keys used above: mutti2025, jaito2026, nortonbaker2025.
% ---------------------------------------------------------------------
% Mutti G, Ocana-Pallares E, Gabaldon T. Newly Developed Structure-Based Methods Do
%   Not Outperform Standard Sequence-Based Methods for Large-Scale Phylogenomics.
%   Molecular Biology and Evolution. 2025;42(7):msaf149. doi:10.1093/molbev/msaf149.
% Norton-Baker B, Komp E, Gado JE, Denton MCR, Mathews II, Murphy NP, Erickson E,
%   Storment OO, Sarangi R, Gauthier NP, McGeehan JE, Beckham GT. Machine
%   Learning-Guided Identification of PET Hydrolases from Natural Diversity. ACS
%   Catalysis. 2025;15(18):16070-16083. doi:10.1021/acscatal.5c03460.
% Jaito N, Kaewsawat N, Sangawthong K, Uengwetwanit T. Identification of novel
%   metagenomic lipases through integrated structural and sequence-based analysis.
%   PeerJ. 2026;14:e20462. doi:10.7717/peerj.20462.
%
% BibTeX (if using a .bib):
% @article{mutti2025, author={Mutti, Giacomo and Oca\~na-Pallar\`es, Eduard and
%   Gabald\'on, Toni}, title={Newly Developed Structure-Based Methods Do Not Outperform
%   Standard Sequence-Based Methods for Large-Scale Phylogenomics},
%   journal={Molecular Biology and Evolution}, year={2025}, volume={42}, number={7},
%   pages={msaf149}, doi={10.1093/molbev/msaf149}}
% @article{nortonbaker2025, author={Norton-Baker, Brenna and Komp, Evan and Gado,
%   Japheth E. and Denton, Mackenzie C. R. and Mathews, Irimpan I. and Murphy, Natasha
%   P. and Erickson, Erika and Storment, Olateju O. and Sarangi, Ritimukta and
%   Gauthier, Nicholas P. and McGeehan, John E. and Beckham, Gregg T.},
%   title={Machine Learning-Guided Identification of PET Hydrolases from Natural
%   Diversity}, journal={ACS Catalysis}, year={2025}, volume={15}, number={18},
%   pages={16070--16083}, doi={10.1021/acscatal.5c03460}}
% @article{jaito2026, author={Jaito, Nongluck and Kaewsawat, Nattha and Sangawthong,
%   Kamollak and Uengwetwanit, Tanaporn}, title={Identification of novel metagenomic
%   lipases through integrated structural and sequence-based analysis}, journal={PeerJ},
%   year={2026}, volume={14}, pages={e20462}, doi={10.7717/peerj.20462}}


% ---------------------------------------------------------------------
% [BLOCK 10]  FIGURE CAPTIONS — updated Fig 4 (left) + new Supplementary Fig.
% ---------------------------------------------------------------------
% Fig 4 caption (left panel) — replace with:
\textbf{The PET-hydrolase branch is not enriched above the enlarged random baseline.}
(Left) Exposed-cleft rate among triad-bearers for the top-300 PET-hydrolase branch
(32.65\%), the acetylcholinesterase branch (8.75\%), and the enlarged random floor
({{base_rate_pct}}\%, $n={{base_n}}$ triad-bearers; Wilson 95\% CI error bar). A two
one-sided test at a pre-registered $\pm${{delta_pp}}-point margin rejects enrichment of
the PET branch above the floor (one-sided 95\% upper bound on the enrichment
{{upper_bound_pp}} points). Regenerated as \texttt{figures/fig4\_left\_enlarged}.

% New Supplementary Figure caption:
\textbf{Supplementary Figure (\label{figS_plddt}). Predicted-confidence is not lower in
the divergent tail.} ECDF (left) and distribution (right) of global pLDDT for the
sub-25\% divergent screened band, its $<$20\% subset, the enlarged random-baseline
triad-bearers, and the near-homolog above-line hits. The divergent band is not the
lower-confidence group (divergent vs.\ baseline Mann--Whitney $p={{plddt_mwu_p}}$,
Cliff's $\delta={{plddt_cliffs}}$), so the exposed-cleft depletion in the divergent
tail is not an ESMFold prediction-error artifact. Source:
\texttt{figures/figS\_plddt\_bands}.
